\documentclass[11pt]{article}
\usepackage[margin=1in]{geometry}
\usepackage{amsmath,amssymb,amsthm}
\usepackage[T1]{fontenc}
\usepackage{lmodern}
\usepackage[colorlinks=true,allcolors=blue,breaklinks=true]{hyperref}
\Urlmuskip=0mu plus 1mu
\hypersetup{pdftitle={Exceptional times for nearest-neighbor exclusion percolation on the triangular lattice have Hausdorff dimension 31/36},pdfauthor={}}
\newtheorem{theorem}{Theorem}
\newtheorem{lemma}[theorem]{Lemma}
\newcommand{\Pp}{\mathbb P}
\newcommand{\E}{\mathbb E}
\newcommand{\T}{\mathbb T}
\newcommand{\R}{\mathbb R}
\newcommand{\one}{\mathbf 1}
\newcommand{\dd}{\,\mathrm d}
\newcommand{\TV}{\mathrm{TV}}
\newcommand{\dimH}{\dim_{\mathrm H}}
\title{Exceptional times for nearest-neighbor exclusion percolation on the triangular lattice have Hausdorff dimension $31/36$}
\author{}
\date{}
\makeatletter
\renewcommand{\maketitle}{\begin{center}{\Large\bfseries\@title\par}\end{center}\vspace{0.5em}}
\makeatother
\begin{document}
\maketitle

\begin{abstract}
We prove the $31/36$ exceptional-time theorem for
nearest-neighbor exclusion percolation on the triangular lattice. The
proof compares two-time crossing probabilities for conservative dimer
updates and independent refresh, using a local reference response to
normalize time. A forest estimate controls every degree of the expansion;
a sewing argument identifies each fixed degree. The resulting annular
mixing estimate gives the energy bound needed for exceptional times.
A Lean formalization accompanies the proof.
\end{abstract}

\section{The question and the statement}

Identify the triangular lattice with
$\T=\{m+n e^{i\pi/3}:m,n\in\mathbb Z\}$, with edges between points at
distance one. Initially, sites are independently black or white with
probability $1/2$. Each unordered edge has an independent rate-one Poisson
clock; when it rings, its two endpoint colors are exchanged. Write
$\omega_t$ for the configuration and
\[
 \mathcal E=\{t\geq0:\omega_t\text{ has an infinite black cluster}\}.
\]
At every deterministic time the configuration is critical Bernoulli
percolation, so it has no infinite cluster almost surely. The question
concerns the random set $\mathcal E$.

Broman, Garban and Steif asked whether conservative percolation has
exceptional times \cite[Section 9.4]{BGS}. Garban and Vanneuville later
formulated the dimension-$31/36$ conjecture for symmetric exclusion
dynamics \cite[Conjecture 1.7]{GV}; it appears in their May 2016 preprint.
Nearest-neighbor exchanges are a principal case of that conjecture.
For independent dynamical percolation the dimension is $31/36$
\cite[Theorem 1.4(1)]{GPSspectral}, while the exclusion upper bound is already known
\cite[Proposition 3.3]{GV}. We prove the nearest-neighbor case on $\T$.

\begin{theorem}\label{thm:main}
Almost surely, simultaneously for all $0\leq a<b<\infty$,
\[
             \dimH\bigl(\mathcal E\cap(a,b)\bigr)=\frac{31}{36}.
\]
In particular, every such interval contains exceptional times.
\end{theorem}

\section{Proof}

\subsection*{1. A local clock and a comparison of finite crossing words}

Put $\varepsilon=2^{-n}$ and
$a_\varepsilon(r)=\Pp(A_4(r,1))$ on the \emph{same} mesh $\varepsilon$.
Here $A_4(r,R)$ denotes four disjoint arms of alternating colors between
the circular cuts at radii $r$ and $R$, using closed-cell paths.
We use RSW, arm separation and quasi-multiplicativity in
the forms of \cite[Sections 2.3 and 4, Propositions 15--16]{Nolin},
with the fixed-shape comparisons of \cite[Section 2]{OAI}.

Fix a sufficiently large microscopic constant $B$. Separate four
alternating arms at the microscopic cut and at a fixed positive radius.
Disjoint RSW corridors attach their outer ends to the left and right
black sides and the upper and lower white sides of
$Q_0=[-1,1]\times[-1/2,1/2]$. A fixed finite collar routes the inner
ends around the neighboring pair, leaving its two bits free; finite
energy pays only a constant factor for its prescribed colors.
Jordan separation \cite[Section 4.1]{Diestel} excludes bypasses across
the white connections, so the crossing reads the first free bit for all
four assignments of the pair. Arm separation and
quasi-multiplicativity therefore give probability at least
$c\,a_\varepsilon(B\varepsilon)$, uniformly over a fixed bulk range.
This is an ordinary one-time attachment estimate.

We also use the ordinary two-source sewing construction of
\cite[Section 6, Lemmas 6.5--6.6 and the relative many-trial estimate]{OAI}.
For a fixed finite polygon family and a compact, positively separated
insertion range, it supplies annular success events invariant under
source exchange and the relevant inside exchange. Success makes the
crossings depend on the enclosed colors through a single join bit.
For two independent ordinary fair fields $\omega,\eta$, the first-point
failure event $N$ satisfies
\begin{equation}\label{eq:seam-input}
 \Pp\bigl(N\cap A_4^{\rm mic}(\omega)\cap A_4^{\rm sq}(\eta)\bigr)
       \leq\epsilon\,a_\varepsilon(B\varepsilon)a_\varepsilon(h).
\end{equation}
The two arm events refer respectively to the microscopic insertion and
the fixed square. To apply the source estimate, trim their paths to
buffered graph cuts inside the protected annulus; the resulting arms
imply the source's probe event. Fixed-shape arm comparisons change its
normalization only by a constant factor. The relative many-trial bound
has the form $\rho+C/n$: choose $\rho$ small, fix the microscopic
enlargement, then choose the finite trial count $n$ large. Next fix $h$
below those trials and finally take a sufficiently fine mesh. This gives
\eqref{eq:seam-input}, uniformly on the compact separated location range.
The reference-square influence is at least
$c\,a_\varepsilon(h)$, by the specialization of
\cite[Lemma 6.1]{OAI}, again using the fixed square/circle comparison.
These estimates apply on the
triangular lattice as stated in \cite[Sections 2 and 6]{OAI}.

Use integer coordinates on $\T$. For each phase $\sigma\in\{0,1\}$,
pair $(2m+\sigma,n)$ with $(2m+\sigma+1,n)$. Consider two rate-one
updates on each pair, independently between pairs:
$C$ replaces its colors uniformly subject to their black count;
$M$ chooses an endpoint uniformly and refreshes that bit fairly.
Thus $C$ keeps or swaps the pair with equal probability, whereas $M$
is independent site refresh at rate $1/2$.

For $A\in\{C,M\}$, let $S_A(\varepsilon)$ be the expected change
indicator of the $Q_0$ crossing when one $A$ update is applied to the
central horizontal pair, with all other colors an ordinary fair field.
Define
\begin{equation}\label{eq:clock}
 J_A(\varepsilon)=\frac{S_A(\varepsilon)}{2\varepsilon^2},
 \qquad u_A(\varepsilon)=\frac{\varepsilon^2}{S_A(\varepsilon)}.
\end{equation}
The attachment estimate and the exact two-bit calculation give
$S_A\geq c\,a_\varepsilon(B\varepsilon)>0$. The four-arm exponent
$5/4$ \cite[Theorem 4]{SW} consequently gives, for every
$0<\eta<3/4$ and all sufficiently small meshes,
\begin{equation}\label{eq:clock-bound}
       u_A(\varepsilon)\leq C_\eta\varepsilon^{3/4-\eta}.
\end{equation}
Translation of the reference rectangle identifies the response
at every insertion with this same $S_A$.

Let $F_\varepsilon$ be a Boolean word in the literal closed-cell crossings
of a fixed finite family of dyadic polygonal quads. Here the quads are
images of dyadic axis-parallel grid quads under
$(u,v)\mapsto u+e^{i\pi/3}v$; invariance of domain
\cite[Theorem 2B.3]{Hatcher} identifies their polygonal interiors. Put
$q^A_\varepsilon(t)=\Pp(F_\varepsilon(\omega_0)=
F_\varepsilon(\omega_t)=1)$ for the corresponding matching dynamics.

\begin{lemma}\label{lem:word}
For each fixed $T>0$ and each phase,
\begin{equation}\label{eq:word}
 q^C_\varepsilon(Tu_C(\varepsilon))-
 q^M_\varepsilon(Tu_M(\varepsilon))\longrightarrow0.
\end{equation}
The second term is the independent site-refresh probability at time
$T\varepsilon^2/(2S_M(\varepsilon))$.
\end{lemma}

\begin{proof}
Take a fixed physical box containing every tested quad and cover its
sites by the matching. The two-time law of one pair has an exact
expansion
\[
 p_A(t)=p_0+t\{w_A+r_A(t)\},\qquad
 \sum w_A=\sum r_A(t)=0,\quad
 \|r_A(t)\|_\TV\leq t\quad(0<t\leq1).
\]
Here $\|v\|_\TV=\sum_h|v(h)|$ is the full variation norm, and
$p_0$ is the fair static pair, constant between the two times;
$w_A$ is the derivative of the actual pair law at zero. Direct
calculation gives $\|w_C\|_\TV=1/2$ and $\|w_M\|_\TV=1$.

For an unordered selection $S$ of pairs, integrate arbitrary zero-mass
signed weights $w_x$ at $S$ and the static law elsewhere. The absolute
value of the resulting kernel is bounded by
\begin{equation}\label{eq:geometry}
          G_\varepsilon(S)\prod_{x\in S}\|w_x\|_\TV,
          \qquad G_\varepsilon(S)\geq0.
\end{equation}
Centering forces a local change at every selected pair. Away from the
quad boundaries this forces four alternating arms; near a side it
forces the appropriate half-plane arms \cite[Lemma 2.2]{OAI}.
Record successive mergers of the probe neighborhoods as a rooted
forest. RSW, arm separation, the four-arm power bounds and the
half-plane three-arm bound \cite[Theorems 20 and 23(ii)]{Nolin}
control the disjoint probe costs.

The two quantitative consequences are as follows. Choose a lower
four-arm power $g\in(5/4,2)$, for example $g=3/2$, and set
$\beta=1-g/2>0$.
For some $A<\infty$, uniformly in phase and update rule, eventually
\emph{for every} $k\geq0$,
\begin{equation}\label{eq:all-degrees}
 u_A(\varepsilon)^k\sum_{|S|=k}G_\varepsilon(S)
       \leq A^k k^{-\beta k},
\end{equation}
where the right side is one at $k=0$. The forest integration uses
the stopping radius $R_k=\min\{R_*,c_*/\sqrt{k}\}$.
Packing ensures that the microscopic cutoff lies below $R_k$ for
every nonempty selection; larger degrees have no selections. Counting
the merge forests and integrating their radial edges gives
\eqref{eq:all-degrees}, including the selection factorial.
Moreover, for fixed $k$ the normalized ordered kernels have a uniform
$L^p$ forest bound for any $p>1$ sufficiently close to one that $gp<2$.

Write $K_A(\varepsilon;\boldsymbol x)$ for the kernel with $k$ raw
insertions $w_A$. On tuples separated from each other and from the
polygon frontiers, the sewing estimate gives
\begin{equation}\label{eq:good}
 \sup_{\boldsymbol x}
 \left|\frac{K_C(\varepsilon;\boldsymbol x)}{S_C(\varepsilon)^k}
       -\frac{K_M(\varepsilon;\boldsymbol x)}{S_M(\varepsilon)^k}\right|
       \longrightarrow0.
\end{equation}
Indeed, hold the other $k-1$ signed insertions fixed and compare the
distinguished insertion and the translated reference crossing with one
common forced square. Put $a=a_\varepsilon(B\varepsilon)$ and
$b=a_\varepsilon(h)$. If $H_h$ is the square kernel and $D_h$ its
reference response, write $\widehat K_A$ for the kernel with distinguished
rule $A$ and the other rules unchanged. The two errors are
$E_A=\widehat K_A D_h-H_h S_A$. The source symmetries and disjoint ordinary
probes give $|E_A|\leq C_k\epsilon a^k b$. The exact cancellation
\[
 D_h(\widehat K_C S_M-\widehat K_M S_C)=S_M E_C-S_C E_M
\]
and the bounds $D_h\geq cb$, $ca\leq S_A\leq Ca$ compare the two
normalized insertions. The extra $k-1$ probe cutoffs are fixed
independently of the tolerance. Telescoping replaces all $k$ entries.
For each tolerance one chooses the finite trials, then a sufficiently
small fixed $h$, and finally sends $\varepsilon$ to zero. The normalized
microscopic difference is independent of $h$, so this proves
\eqref{eq:good}.

The degree-$k$ coefficient after clock normalization is exactly
\begin{equation}\label{eq:coefficient}
 c_{A,k}(\varepsilon)=\frac{\varepsilon^{2k}}{k!}
  \sum_{\boldsymbol x\ \mathrm{distinct}}
       \frac{K_A(\varepsilon;\boldsymbol x)}{S_A(\varepsilon)^k}.
\end{equation}
Remove tuples within distance $\delta$ of a collision or a polygon
frontier. Their scaled counting mass tends to zero with $\delta$;
the $L^p$ forest bound and H\"older's inequality make their contribution
uniformly negligible. On the remaining tuples use \eqref{eq:good}.
Thus $c_{C,k}-c_{M,k}\to0$ for each fixed $k$, including the common
static degree zero.

Finally, \eqref{eq:geometry} also applies to the centered remainders
$r_A$. By \eqref{eq:clock-bound}, replacing raw insertions by the
effective insertions changes each fixed-degree coefficient by $o(1)$.
Equation \eqref{eq:all-degrees} bounds the degree-$k$ terms at time $T$
by $(2AT)^k k^{-\beta k}$, a summable sequence. Dominated summation
therefore proves \eqref{eq:word} for the actual finite product laws,
including all repeated-update remainders. The identification of $M$
with site refresh at half its pair time is exact.
\end{proof}

\subsection*{2. From crossing words to an energy bound}

The dense quad coordinates and arm-event convergence of
\cite[Definition 2.1, Lemma 2.9 and Corollary 2.10]{GPSpivotal}
approximate a fixed hexagonal four-arm event in static probability
by a finite crossing word: first approximate its limiting indicator by
a finite cylinder, then pass to sufficiently fine meshes. The word is
fixed before the mesh threshold. Stationarity bounds the error in a
two-time probability by twice the static error. Apply
Lemma~\ref{lem:word}, followed by full independent mixing
\cite[Theorems 1.3(i), 1.4(i) and Remark 1.5, $\gamma=0$]{GHSS}.
The local independent clock is admissible because
$S_M(\varepsilon)\leq C\alpha_4(\varepsilon,1)$ by ordinary pivotal
necessity; hence its site time is at least a fixed multiple of
$T\varepsilon^2/\alpha_4(\varepsilon,1)$. Same-event independent
correlations decrease with time by their Fourier expansion, placing
the local clock in the mixing regime.

The square of the static four-arm probability of an annulus of aspect
$M$ is at most $CM^{-9/4}$. Choose a sufficiently large dyadic $M$,
then a sufficiently large fixed $T$, then a sufficiently fine mesh.
The resulting conservative four-arm bound is at most
$(2M)^{-17/8}$. Its reference time at lattice scale $s$ is bounded by
$C_\eta s^{-(3/4-\eta)}$, by \eqref{eq:clock-bound}. Buffered disjoint
annuli and monotonicity in time iterate this fixed-aspect estimate.
Dyadic rounding produces a decorrelation length
\begin{equation}\label{eq:scale}
       m(t)\leq C_\eta t^{-1/(3/4-\eta)},\qquad 0<t<1,
\end{equation}
above which the paired four-arm probes decay with exponent $17/8>2$.

Let $A_R=\{0\leftrightarrow\partial B_R\}$ and
$\alpha_1(R)=\Pp(A_R)$. For canonical matching dynamics write
$q_R^C(t)=\Pp(A_R(\omega_0)\cap A_R(\omega_t))$.
For a wired annular connection from $r$ to $R$, average the exterior
replicas independently and retain canonical dependence only in a
buffered middle annulus. Write $q_{\rm mix}$ for this mixed two-time
probability and $a(r,R)$ for its ordinary one-time mean. Finite-cluster
boundary separation is used in the form of \cite[Section 5.5]{Grimmett}.
The exact
conditional independent-replica representation permits the finite
Efron--Stein inequality \cite[Proposition 1]{BBLM}. A changed conditional connection
probability forces a paired four-arm probe. All active cells lie in the
buffered bulk, and the exponent $17/8=2+1/8$ gives
\[
 q_{\rm mix}-a(r,R)^2
       \leq C\bigl(m(t)/r\bigr)^{1/8}q_{\rm mix}.
\]
The duplicated-event Harris inequality and fixed black collars give
$q^C(W(r,R))\leq c^{-8}q_{\rm mix}$ for the canonical wired event.
Choose $N$ so that the variance coefficient is at most $1/2$ when
$r\geq Nm(t)$. Absorption at $r=2Nm(t)$ bounds this canonical
probability by $2c^{-8}a(r,R)^2$. The disjoint inner and outer matching
labels and ordinary one-arm quasi-multiplicativity then yield
\[
 q_R^C(t)\leq C\,
       \frac{\alpha_1(R)^2}{\alpha_1(Nm(t))}
\]
in the large-radius case; the static bound handles smaller radii.
The one-arm exponent and quasi-multiplicativity \cite{LSW,Nolin},
together with \eqref{eq:scale}, give, for every $\zeta>0$,
\begin{equation}\label{eq:onearm}
 q_R^C(t)\leq
 \min\bigl\{\alpha_1(R),\,
             C_\zeta\alpha_1(R)^2t^{-5/36-\zeta}\bigr\}.
\end{equation}
The numerical exponent is
$(5/48)/(3/4)=5/36$. All finite covers and both matching phases are
included in these estimates.

\subsection*{3. Nearest-neighbor exclusion and exceptional times}

Let $H_C$ be the positive generator of canonical matching dynamics
and $H_X$ that of rate-one nearest-neighbor swaps, on a finite edge
truncation containing the event's matching-closed support. On a matched edge,
\[
 I-\Pi=\tfrac12(I-\mathsf{Swap}),
 \qquad H_X\geq H_C\geq0
\]
as quadratic forms. For $0<\gamma<31/36$, choose $\zeta>0$ with
$\gamma+5/36+\zeta<1$. Equation \eqref{eq:onearm} gives
\[
 \sup_R\alpha_1(R)^{-2}
       \int_0^1t^{-\gamma}q_R^C(t)\dd t<\infty.
\]
Fractional resolvent order transfers this bound to exclusion. More
explicitly, $(I+H_X)^{\gamma-1}\leq(I+H_C)^{\gamma-1}$; the Laplace
representation, followed by monotonicity of the canonical correlation
for the tail $t\geq1$, gives
\begin{equation}\label{eq:energy}
 \int_0^1t^{-\gamma}q_R^X(t)\dd t
       \leq(e+1)\int_0^1t^{-\gamma}q_R^C(t)\dd t
       \leq C_\gamma\alpha_1(R)^2.
\end{equation}
Backward-trace stabilization, finite-cylinder convergence
and Fatou's lemma pass the bound along an edge exhaustion to the
actual infinite process; the containment threshold may depend on $R$.

On $I=[0,1]$, consider the normalized occupation
measures
\[
       \nu_R(\dd t)=
       \alpha_1(R)^{-1}\one_{A_R(\omega_t)}\one_I(t)\dd t.
\]
Their expected mass is one. Stationarity and \eqref{eq:energy} bound
their expected $\gamma$-energy uniformly in $R$; energy also bounds
squared mass on $I$. A second-moment and Markov argument therefore
gives, with positive probability, a common positive mass lower bound
and finite energy upper bound for infinitely many radii. Apply the
compact-set capacity criterion \cite[Lemma 6.2]{SS} to the decreasing
closures of the arm-time sets. Their intersection has dimension at
least $\gamma$ on this measurable witness event.

Closing each site's black-time set gives a configuration with an infinite
cluster at every time in this compact intersection. There are almost
surely no simultaneous rings of distinct edge clocks. Local path
regularity shows that this configuration differs from the actual one
at at most the two endpoints of the edge ringing at that time.
Finite recoloring preserves existence of some infinite
cluster on a locally finite graph. Thus the compact intersection consists
of actual exceptional times, although its clusters need not contain
the origin.

The product construction from Poisson marks \cite{LP} and graphical
regularity \cite[Appendix A]{GV} is equivariant under lattice translation
and has an ergodic unit spatial shift. Saturating the measurable
positive-probability witness event under this shift upgrades the lower
bound to probability one. Affine pullback to a time window $a+c[0,1]$,
$0<c\leq1$, changes the energy bound by at most $c^{\gamma-1}$.
Intersect over rational $\gamma\in(0,31/36)$ and rational windows;
inclusion of rational subwindows gives the simultaneous lower bound
for every real $0\leq a<b$.
The upper bound follows from \cite[Proposition 3.3]{GV}, after the proved
constant time change from its normalized rate convention. Hausdorff
dimension is invariant under that positive linear time change. This
proves Theorem~\ref{thm:main}.

\paragraph{Formalization.}
The accompanying Lean 4.34.1 development verifies the argument from
24 explicit formulations of the cited results, together with positivity
of the one-arm decay exponent. The literature results and their
specializations explained above are supplied as hypotheses; the new
comparison, coefficient limits, forest bounds, clock estimates and
exceptional-time argument are proved within Lean. The entry point is
\texttt{exists\_realization\_dimension} in the module
\texttt{PublishedNearestNeighborDimension}. The complete build checks
443 source files, with only \texttt{propext}, \texttt{Classical.choice}
and \texttt{Quot.sound} as transitive axioms. The accompanying archive
\texttt{nearest-neighbor-31-36-lean.zip} contains the precise input
formulations, source correspondences and reproducible build record.

\begin{thebibliography}{99}
\small
\bibitem{BGS}
E. I. Broman, C. Garban and J. E. Steif,
\emph{Exclusion sensitivity of Boolean functions},
Probab. Theory Relat. Fields \textbf{155} (2013), 621--663.
\href{https://arxiv.org/abs/1101.1865v1}{arXiv:1101.1865v1} (2011).

\bibitem{GV}
C. Garban and H. Vanneuville,
\emph{Exceptional times for percolation under exclusion dynamics},
Ann. Sci. \'Ec. Norm. Sup\'er. (4) \textbf{52} (2019), 1--57.
\href{https://arxiv.org/abs/1605.04766v3}{arXiv:1605.04766v3};
\href{https://doi.org/10.24033/asens.2383}{doi:10.24033/asens.2383}.

\bibitem{GPSspectral}
C. Garban, G. Pete and O. Schramm,
\emph{The Fourier spectrum of critical percolation},
Acta Math. \textbf{205} (2010), 19--104.
\href{https://arxiv.org/abs/0803.3750v4}{arXiv:0803.3750v4}
(corrected version, 2013).

\bibitem{Nolin}
P. Nolin,
\emph{Near-critical percolation in two dimensions},
Electron. J. Probab. \textbf{13} (2008), 1562--1623.
\href{https://arxiv.org/abs/0711.4948v1}{arXiv:0711.4948v1}:
Propositions 15--16 and Theorems 20 and 23.

\bibitem{LSW}
G. F. Lawler, O. Schramm and W. Werner,
\emph{One-arm exponent for critical 2D percolation},
Electron. J. Probab. \textbf{7} (2002), paper no. 2.
\href{https://arxiv.org/abs/math/0108211}{arXiv:math/0108211}, Theorem 1.1.

\bibitem{SW}
S. Smirnov and W. Werner,
\emph{Critical exponents for two-dimensional percolation},
Math. Res. Lett. \textbf{8} (2001), 729--744.
\href{https://arxiv.org/abs/math/0109120v2}{arXiv:math/0109120v2}, Theorem 4.

\bibitem{GPSpivotal}
C. Garban, G. Pete and O. Schramm,
\emph{Pivotal, cluster and interface measures for critical planar percolation},
J. Amer. Math. Soc. \textbf{26} (2013), 939--1024.
\href{https://arxiv.org/abs/1008.1378v5}{arXiv:1008.1378v5},
Lemma 2.9 and Corollary 2.10.

\bibitem{GHSS}
C. Garban, N. Holden, A. Sep\'ulveda and X. Sun,
\emph{Liouville dynamical percolation},
Probab. Theory Relat. Fields \textbf{180} (2021), 621--678.
\href{https://arxiv.org/abs/1905.06940v2}{arXiv:1905.06940v2},
Theorems 1.3(i), 1.4(i) and Remark 1.5.

\bibitem{OAI}
OpenAI,
\emph{From critical crossings to quenched near-critical universality
in Voronoi percolation}, preprint, 5 October 2026.
\href{https://github.com/openai/math/tree/adc7f1241b42e322a6451854ab7e4b4c146bf78a/preprints/From-critical-crossings-to-quenched-near-critical-universality-in-Voronoi-percolation-October-5-2026/build/sections}{Pinned source files},
commit \texttt{adc7f1241b42e322a6451854ab7e4b4c146bf78a},
Sections 2 and 6.

\bibitem{BBLM}
S. Boucheron, O. Bousquet, G. Lugosi and P. Massart,
\emph{Moment inequalities for functions of independent random variables},
Ann. Probab. \textbf{33} (2005), 514--560.
\href{https://arxiv.org/abs/math/0503651}{arXiv:math/0503651}, Proposition 1.

\bibitem{SS}
O. Schramm and J. E. Steif,
\emph{Quantitative noise sensitivity and exceptional times for percolation},
Ann. of Math. (2) \textbf{171} (2010), 619--672.
\href{https://annals.math.princeton.edu/wp-content/uploads/annals-v171-n2-p01-p.pdf}{Published text},
Lemma 6.2, p. 650.

\bibitem{LP}
G. Last and M. Penrose,
\emph{Lectures on the Poisson Process}, Cambridge University Press, 2017.
\href{https://stoch.math.kit.edu/img/Last/lastpenrose2017.pdf}{Author manuscript},
Definition 3.1, Theorem 3.6 and Proposition 6.9.

\bibitem{Diestel}
R. Diestel, \emph{Graph Theory}, 5th ed., Springer, 2017, Section 4.1.

\bibitem{Grimmett}
G. Grimmett, \emph{Probability on Graphs}, Cambridge University Press, 2010.
\href{https://www.statslab.cam.ac.uk/~grg/books/USpgs-rev3.pdf}{Author's 2012 revision},
Section 5.5.

\bibitem{Hatcher}
A. Hatcher, \emph{Algebraic Topology}, Cambridge University Press, 2002.
\href{https://pi.math.cornell.edu/~hatcher/AT/ATch2.pdf}{Chapter 2}, Theorem 2B.3.
\end{thebibliography}
\end{document}
